
\section{\BookRecordIndex{T0275}气泡相连}\label{sec:4-7}
\zhindex{qipaoxianglian}{气泡相连}

本节始终假设 $(M,\omega)$ 是紧\BookRecordIndex{T0358}辛流形，且
$J\in\mathcal J_\tau(M,\omega)$ 是 $M$ 上被 $\omega$ \BookRecordIndex{T0372}驯服的\BookRecordIndex{T0088}殆复结构。我们只讨论
$J$-全纯球面的内部\BookRecordIndex{T0243}冒泡。下面命题~\ref{prop:4-7-1} 和命题~\ref{prop:4-7-2}
所陈述的主要结果是：适当选择\BookRecordIndex{T0412}重标度序列 $\psi^\nu$ 后，\BookRecordIndex{T0275}气泡 $v$ 会与原映射列
$u^\nu$ 的极限 $u$ 相连。证明的关键是选择\BookRecordIndex{T0412}重标度序列，使能量没有损失；见下面的
恒等式 (iii)。稳定性条件 (iv) 的重要性将在下一章显现。

在下面的命题中，我们把 $S^2$ 等同于黎曼球面 $\mathbb C\cup\{\infty\}$，并以
$G=\operatorname{PSL}(2,\mathbb C)$ 表示 \BookRecordIndex{T0040}Möbius 变换群。

\begin{proposition}\label{prop:4-7-1}
设 $J^\nu\in\mathcal J_\tau(M,\omega)$ 是 $M$ 上一列被 $\omega$ \BookRecordIndex{T0372}驯服的\BookRecordIndex{T0088}殆复结构，
并在 $C^\infty$ 拓扑下收敛到 $J$。固定一点 $z_0\in\mathbb C$ 和一个数 $r>0$。
设 $u^\nu:B_r(z_0)\to M$ 是一列 $J^\nu$-全纯曲线，
$u:B_r(z_0)\to M$ 是一条 \BookRecordIndex{T0003}$J$-全纯曲线，并且满足：
\begin{enumerate}[label=(\alph*),leftmargin=3em]
\item $u^\nu$ 在 $B_r(z_0)\setminus\{z_0\}$ 的每个紧子集上以 $C^\infty$ 收敛到 $u$；
\item 极限
\[
  m_0:=\lim_{\varepsilon\to0}\lim_{\nu\to\infty}
  E(u^\nu;B_\varepsilon(z_0))
\]
存在且为正。
\end{enumerate}
则存在一个子列（仍记为 $u^\nu$）、一列 \BookRecordIndex{T0040}Möbius 变换 $\psi^\nu\in G$、一条
$J$-全纯球面 $v:S^2\to M$，以及有限多个两两不同的点
$z_1,\ldots,z_\ell,z_\infty\in S^2$，使得：
\begin{enumerate}[label=(\roman*),leftmargin=3em]
\item $\psi^\nu$ 在 $S^2\setminus\{z_\infty\}$ 的每个紧子集上以 $C^\infty$
收敛到 $z_0$；
\item 映射列 $v^\nu:=u^\nu\circ\psi^\nu$ 在
$S^2\setminus\{z_1,\ldots,z_\ell,z_\infty\}$ 的每个紧子集上以 $C^\infty$
收敛到 $v$，并且对 $j=1,\ldots,\ell$，极限
\[
  m_j:=\lim_{\varepsilon\to0}\lim_{\nu\to\infty}
  E(v^\nu;B_\varepsilon(z_j))
\]
存在且为正；
\item $E(v)+\sum_{j=1}^{\ell}m_j=m_0$；
\item 若 $v$ 是常值映射，则 $\ell\ge2$。
\end{enumerate}
\end{proposition}

下一个命题断言：$u$ 的像与 $v$ 的像在点 $u(z_0)=v(z_\infty)$ 处相连。这是因为上述
条件 (iii) 表明，两者之间没有余下的能量可形成新的\BookRecordIndex{T0275}气泡。

\begin{proposition}\label{prop:4-7-2}
设 $J^\nu\in\mathcal J_\tau(M,\omega)$、$z_0\in\mathbb C$ 以及
$u,u^\nu:B_r(z_0)\to M$ 如命题~\ref{prop:4-7-1} 的假设，并设
$\psi^\nu\in G$、$v:S^2\to M$ 以及
$z_1,\ldots,z_\ell,z_\infty\in S^2$ 满足命题~\ref{prop:4-7-1} 的 (i)--(iii)。则
\[
  u(z_0)=v(z_\infty).
\]
此外，对每个 $\varepsilon>0$，都存在常数 $\delta>0$ 和
$\nu_0\in\mathbb N$，使得对每个整数 $\nu\ge\nu_0$ 和每个 $z\in S^2$，
\[
  d(z,z_0)+d\bigl((\psi^\nu)^{-1}(z),z_\infty\bigr)<\delta
  \quad\Longrightarrow\quad
  d\bigl(u^\nu(z),u(z_0)\bigr)<\varepsilon.
\]
\end{proposition}

命题~\ref{prop:4-7-1} 和命题~\ref{prop:4-7-2} 的证明，都以如下关于任意长圆柱上
\BookRecordIndex{T0003}$J$-全纯曲线能量的结果为基础。它断言：若能量充分小，则能量不可能均匀地铺开，
而必定集中在两端附近。我们用闭\BookRecordIndex{T0170}环域来表述这个结果：
\[
  A(r,R):=B_R\setminus\operatorname{int}(B_r),
  \qquad r<R.
\]
下面半径 $r,R$ 的相对大小至关重要。应当把比值 $R/r$ 想成非常大，并且远大于
$e^T$。在下面的式~\eqref{eq:4-7-2} 中，$d$ 表示由 $\omega$ 和 $J$ 所给黎曼度量
$g_J$ 诱导的 $M$ 上距离函数。

\begin{lemma}\label{lem:4-7-3}
设 $(M,\omega)$ 是紧\BookRecordIndex{T0358}辛流形，且 $J\in\mathcal J_\tau(M,\omega)$。则存在常数
$0<\mu<1$、$\delta>0$ 和 $c>0$，使得下述结论成立。若 $0<r<R$，且
$u:A(r,R)\to M$ 是满足能量界
\[
  E(u)=E(u;A(r,R))<\delta
\]
的 \BookRecordIndex{T0003}$J$-全纯曲线，则
\begin{equation}\label{eq:4-7-1}
  E\bigl(u;A(e^Tr,e^{-T}R)\bigr)
  \le\frac{4^\mu}{e^{2\mu T}}E(u)
\end{equation}
以及
\begin{equation}\label{eq:4-7-2}
  \sup_{z,z'\in A(e^Tr,e^{-T}R)}d\bigl(u(z),u(z')\bigr)
  \le\frac c{e^{\mu T}}\sqrt{E(u)}
\end{equation}
对
\[
  \log2\le T\le\log\sqrt{R/r}
\]
成立。
\end{lemma}

\begin{remark}\label{rem:4-7-4}
引理~\ref{lem:4-7-3} 的证明推迟到本节末尾。该证明表明，常数 $\mu,\delta,c$
显式依赖引理~\ref{lem:4-3-1}(i) 和定理~\ref{thm:4-4-1}(iii) 中的常数。因此，由
注~\ref{rem:4-3-2} 和注~\ref{rem:4-4-2}，可以使 $\mu,\delta,c$ 关于
$C^2$ 拓扑连续地依赖于 $J$。
\end{remark}

作为引理~\ref{lem:4-7-3} 的推论，我们得到下面关于 $J^\nu$-全纯\BookRecordIndex{T0170}环域列的结果：
它们无能量损失地收敛到一对 $J$-全纯圆盘。这个结果是命题~\ref{prop:4-7-2}
证明中的主要步骤。仍把 $S^2$ 等同于黎曼球面 $\mathbb C\cup\{\infty\}$。

\begin{lemma}\label{lem:4-7-5}
设 $J^\nu\in\mathcal J_\tau(M,\omega)$ 是 $M$ 上一列被 $\omega$ \BookRecordIndex{T0372}驯服的\BookRecordIndex{T0088}殆复结构，
并在 $C^\infty$ 拓扑下收敛到 $J$。固定正实数 $r$ 和趋于零的正数列
$\delta^\nu$。设 $u^\nu:A(\delta^\nu/r,r)\to M$ 是一列 $J^\nu$-全纯曲线，
满足：
\begin{itemize}
\item $u^\nu$ 在 $B_r\setminus\{0\}$ 的紧子集上一致收敛到一条
\BookRecordIndex{T0003}$J$-全纯曲线 $u:B_r\to M$；
\item $u^\nu(\delta^\nu\,\cdot)$ 在
$\mathbb C\setminus\operatorname{int}B_{1/r}$ 的紧子集上一致收敛到一条
\BookRecordIndex{T0003}$J$-全纯曲线
$v:S^2\setminus\operatorname{int}B_{1/r}\to M$；
\item
\[
  \lim_{\rho\to0}\lim_{\nu\to\infty}
  E\bigl(u^\nu;A(\delta^\nu/\rho,\rho)\bigr)=0.
\]
\end{itemize}
则
\[
  u(0)=v(\infty),
  \qquad
  \lim_{\rho\to0}\lim_{\nu\to\infty}
  \sup_{\delta^\nu/\rho\le|z|\le\rho}
  d\bigl(u^\nu(z),u(0)\bigr)=0.
\]
\end{lemma}

\begin{proof}
简记
\[
  E^\nu(\rho):=E\bigl(u^\nu;A(\delta^\nu/\rho,\rho)\bigr),
  \qquad
  E(\rho):=\lim_{\nu\to\infty}E^\nu(\rho).
\]
则 $E(\rho)$ 在 $\rho\to0$ 时趋于零。因此，只要先取充分小的 $\rho>0$，再取充分大的
$\nu$，就可以把引理~\ref{lem:4-7-3} 应用于
$J^\nu$-全纯曲线 $u^\nu:A(\delta^\nu/\rho,\rho)\to M$。若 $2\rho\le r$，则
\[
  A(\delta^\nu/(2\rho),2\rho)\subset A(\delta^\nu/r,r).
\]
对\BookRecordIndex{T0170}环域 $A(\delta^\nu/(2\rho),2\rho)$ 应用式~\eqref{eq:4-7-2}，并取
$T=\log2$，便得到正数 $\delta,c$，使得
\begin{equation}\label{eq:4-7-3}
  E^\nu(2\rho)<\delta
  \quad\Longrightarrow\quad
  \sup_{z,z'\in A(\delta^\nu/\rho,\rho)}
  d\bigl(u^\nu(z),u^\nu(z')\bigr)
  <c\sqrt{E^\nu(2\rho)}.
\end{equation}
令 $\nu\to\infty$，得到
\[
  E(2\rho)<\delta
  \quad\Longrightarrow\quad
  \begin{aligned}[t]
  d\bigl(u(\rho),v(1/\rho)\bigr)
  &=\lim_{\nu\to\infty}
    d\bigl(u^\nu(\rho),u^\nu(\delta^\nu/\rho)\bigr)\\
  &\le c\sqrt{E(2\rho)}.
  \end{aligned}
\]
再令 $\rho\to0$，便得到 $u(0)=v(\infty)=:x$。

为证明第二个结论，固定 $\varepsilon>0$。我们必须证明存在 $\rho>0$ 和整数
$\nu_0>0$，使得
\[
  \nu\ge\nu_0,
  \quad\frac{\delta^\nu}{\rho}\le|z|\le\rho
  \quad\Longrightarrow\quad
  d\bigl(u^\nu(z),u(0)\bigr)<\varepsilon.
\]
首先选取正数 $\delta,c$，使式~\eqref{eq:4-7-3} 对 $0<\rho\le r/2$ 成立。
其次取充分小的 $\rho>0$，使 $2\rho\le r$ 且
\[
  E(2\rho)<\delta,
  \qquad c\sqrt{E(2\rho)}<\varepsilon/3,
  \qquad d\bigl(u(\rho),u(0)\bigr)<\varepsilon/3.
\]
最后取充分大的 $\nu_0$，使得当 $\nu\ge\nu_0$ 时，
\[
  E^\nu(2\rho)<\delta,
  \qquad c\sqrt{E^\nu(2\rho)}<\varepsilon/3,
  \qquad d\bigl(u^\nu(\rho),u(\rho)\bigr)<\varepsilon/3.
\]
现在设 $\nu\ge\nu_0$ 且 $\delta^\nu/\rho\le|z|\le\rho$。由三角不等式和
式~\eqref{eq:4-7-3}，
\[
  \begin{aligned}
  d\bigl(u^\nu(z),u(0)\bigr)
  &\le d\bigl(u^\nu(z),u^\nu(\rho)\bigr)
    +d\bigl(u^\nu(\rho),u(\rho)\bigr)
    +d\bigl(u(\rho),u(0)\bigr)\\
  &\le c\sqrt{E^\nu(2\rho)}
    +d\bigl(u^\nu(\rho),u(\rho)\bigr)
    +d\bigl(u(\rho),u(0)\bigr)\\
  &<\varepsilon/3+\varepsilon/3+\varepsilon/3=\varepsilon.
  \end{aligned}
\]
这就证明了引理~\ref{lem:4-7-5}。
\end{proof}


\begin{proof}[命题~\ref{prop:4-7-2} 的证明]
在这个命题中，\BookRecordIndex{T0412}重标度序列 $\psi^\nu$ 是任意的。我们将由引理~\ref{lem:4-7-5}
\BookRecordIndex{T0326}推出结论。为此，必须修改\BookRecordIndex{T0412}重标度序列 $\psi^\nu$，使它具有
$\psi^\nu(z)=\delta^\nu z$ 的形式，其中 $\delta^\nu\to0$；然后再验证
引理~\ref{lem:4-7-5} 的假设。这分三步完成。

\medskip
\noindent\textbf{步骤 1.}\quad
\emph{不失一般性，可以假设 $z_0=0$ 且 $z_\infty=\infty$。}

为此，选取等距变换 $\phi_0,\phi_\infty\in G$，使得
$\phi_0(z_0)=0$ 且 $\phi_\infty(z_\infty)=\infty$。然后把
$u^\nu,v^\nu,\psi^\nu$ 以及 $u,v$ 分别替换为
\[
  \widetilde u^\nu=u^\nu\circ\phi_0^{-1},
  \qquad
  \widetilde v^\nu=v^\nu\circ\phi_\infty^{-1},
  \qquad
  \widetilde\psi^\nu
  =\phi_0\circ\psi^\nu\circ\phi_\infty^{-1},
\]
以及
\[
  \widetilde u=u\circ\phi_0^{-1},
  \qquad
  \widetilde v=v\circ\phi_\infty^{-1}.
\]
这些函数满足本命题的假设
\sourcefootnote{原书此处写作“满足该引理的假设”（the assumptions of the lemma）；
但引理~\ref{lem:4-7-5} 所需的重标度形式和无能量损失条件分别到步骤 2、3 才建立，
此处坐标变换保持的是命题~\ref{prop:4-7-2} 的假设，故改正引用对象。}；若它们满足
结论，原来的函数也满足结论。这就证明了步骤 1。

\medskip
\noindent\textbf{步骤 2.}\quad
\emph{不失一般性，可以假设 $\psi^\nu(z)=\delta^\nu z$，其中
$\delta^\nu>0$ 且 $\delta^\nu$ 收敛到零。}

由于 $\psi^\nu(0)\to0$ 且 $(\psi^\nu)^{-1}(\infty)\to\infty$，存在序列
$\rho_0^\nu,\rho_\infty^\nu\in G$，它们分别一致收敛到以某个
$\mathbb C^*$ 中元素作乘法的映射，并且当 $\nu$ 充分大时满足
\[
  \rho_0^\nu\bigl(\psi^\nu(0)\bigr)=0,
  \qquad \rho_0^\nu(\infty)=\infty,
  \qquad
  \rho_\infty^\nu\bigl((\psi^\nu)^{-1}(\infty)\bigr)=\infty,
  \qquad \rho_\infty^\nu(0)=0.
\]
取子列后，还可以假设
\[
  \rho_\infty^\nu(1)=1,
  \qquad
  \rho_0^\nu\bigl(\psi^\nu(1)\bigr)\in\mathbb R^+.
\]
现在把 $u^\nu,v^\nu,\psi^\nu$ 分别替换为
\[
  \widetilde u^\nu:=u^\nu\circ(\rho_0^\nu)^{-1},
  \qquad
  \widetilde v^\nu:=v^\nu\circ(\rho_\infty^\nu)^{-1},
  \qquad
  \widetilde\psi^\nu
  :=\rho_0^\nu\circ\psi^\nu\circ(\rho_\infty^\nu)^{-1}.
\]
这些函数满足步骤 2 的要求。

\medskip
\noindent\textbf{步骤 3.}\quad
\emph{当 $\rho\to0$ 时，
\[
  E(\rho):=\lim_{\nu\to\infty}
  E\bigl(u^\nu;B_\rho\setminus B_{\delta^\nu/\rho}\bigr)
\]
收敛到零。}

由能量的共形不变性，
\[
  \begin{aligned}
  E(\rho)
  &=\lim_{\nu\to\infty}E(u^\nu;B_\rho)
    -\lim_{\nu\to\infty}E(v^\nu;B_{1/\rho})\\
  &=m_0+E(u;B_\rho)-E(v;B_{1/\rho})
    -\sum_{j=1}^{\ell}m_j\\
  &=E(u;B_\rho)+E(v;\mathbb C\setminus B_{1/\rho}).
  \end{aligned}
\]
这里用到了假设中 $u^\nu,v^\nu$ 满足命题~\ref{prop:4-7-1} 的 (a)、(b) 与
(i)--(iii)。上面的第二个等号使用了 (b) 和 (ii)，最后一个等号使用了 (iii)。
步骤 3 得证。

由步骤 3，映射列 $u^\nu$ 满足引理~\ref{lem:4-7-5} 的假设，因而
$u(0)=v(\infty)$。此外，由于 $\psi^\nu(z)=\delta^\nu z$，条件
$\delta^\nu/\rho<|z|<\rho$ 转化为
\[
  |z|<\rho,
  \qquad
  |(\psi^\nu)^{-1}(z)|>1/\rho,
\]
所以它表示 $z$ 接近 $z_0=0$，且 $(\psi^\nu)^{-1}(z)$ 接近
$z_\infty=\infty$。因此，命题~\ref{prop:4-7-2} 的第二个结论也由
引理~\ref{lem:4-7-5} \BookRecordIndex{T0326}推出。命题~\ref{prop:4-7-2} 得证。
\end{proof}

\begin{proof}[命题~\ref{prop:4-7-1} 的证明]
证明分四步。

\medskip
\noindent\textbf{步骤 1.}\quad
\emph{不失一般性，可以假设 $z_0=0$，并且每个函数
$z\mapsto|du^\nu(z)|$ 都在 $z=0$ 处取得最大值。}

由 (a)，对每个 $\varepsilon>0$，序列 $|du^\nu|$ 在\BookRecordIndex{T0170}环域
$B_r(z_0)\setminus B_\varepsilon(z_0)$ 上一致有界；而由 (b)，
\[
  \lim_{\nu\to\infty}\sup_{B_\varepsilon(z_0)}|du^\nu|=\infty
\]
对每个 $\varepsilon>0$ 都成立。因此，任意一列满足
\[
  z^\nu\in B_r(z_0),
  \qquad
  |du^\nu(z^\nu)|=\sup_{B_r(z_0)}|du^\nu|
\]
的点都收敛到 $z_0$。于是可以假设对所有 $\nu$ 都有
$z^\nu\in B_{r/2}(z_0)$。考虑映射
\[
  \widetilde u^\nu:B_{r/2}\to M,
  \qquad
  \widetilde u^\nu(z):=u^\nu(z+z^\nu).
\]
它们满足本命题的假设 (a)、(b)，其中 $r$ 换成 $r/2$，$z_0$ 换成零。此外，若它们对
某个序列 $\psi^\nu$ 满足结论，那么原来的函数 $u^\nu$ 也满足结论，只须把
$\psi^\nu$ 换成 $z\mapsto\psi^\nu(z)+z^\nu$。因此，把 $u^\nu$ 换成
$\widetilde u^\nu$，并把 $r$ 换成 $r/2$，便得到步骤 1。

\medskip
\noindent\textbf{步骤 2.}\quad
\emph{假设步骤 1 成立。令 $\hbar$ 为命题~\ref{prop:4-1-4} 中的常数。由
注~\ref{rem:4-7-4}，可以选取固定的 $0<\mu<1$ 和正数
$0<\delta<\hbar$，使引理~\ref{lem:4-7-3} 的结论在 $J$ 换成 $J^\nu$ 时，对所有
充分大的 $\nu$ 以这些常数一致成立
\sourcefootnote{原书此处固定 $\mu=1/2$；但定理~\ref{thm:4-4-1}(iii) 在一般\BookRecordIndex{T0372}驯服情形
只给某个未受控制的等周常数，而引理~\ref{lem:4-7-3} 证明中的 $\mu$ 必须据此选取，
故不能一般地指定为 $1/2$。这里沿用该引理及注~\ref{rem:4-7-4} 给出的某个固定
$0<\mu<1$；原书步骤 3 的 $\delta/R$ 是 $\mu=1/2$ 时的特例，一般情形的对应项仍趋于零，
所以反证不变。}。于是，当 $\nu$ 充分大时，存在唯一的常数
$\delta^\nu>0$，使得}
\begin{equation}\label{eq:4-7-4}
  E(u^\nu;B_{\delta^\nu})=m_0-\frac\delta2.
\end{equation}
\emph{定义 $\psi^\nu\in G$ 为 $\psi^\nu(z):=\delta^\nu z$。则
$\delta^\nu\to0$，所以 $\psi^\nu$ 满足 (i)，其中 $z_\infty=\infty$。此外，存在一个
子列（仍记为 $u^\nu,\psi^\nu$）满足 (ii)。}

由定理~\ref{thm:4-6-1}，有 $m_0\ge\hbar$。因此，根据 $m_0$ 的定义，当 $\nu$
充分大时，存在唯一的 $\delta^\nu>0$ 满足式~\eqref{eq:4-7-4}；而且
$\delta^\nu$ 收敛到零。我们这样选择 $\delta^\nu$，是为了使集中在零点的“\BookRecordIndex{T0275}气泡”能量
$m_0$ 几乎全部落在球 $B_{\delta^\nu}$ 内。因此，当 $\varepsilon>0$ 充分小时，
\BookRecordIndex{T0170}环域 $B_\varepsilon\setminus B_{\delta^\nu}$ 中没有足够的能量产生另一个\BookRecordIndex{T0275}气泡。

由于 $\delta^\nu\to0$，序列 $\psi^\nu$ 满足命题~\ref{prop:4-7-1} 的条件 (i)，
其中 $z_\infty:=\infty$。现在定义
\[
  v^\nu:B_{r/\delta^\nu}\to M,
  \qquad
  v^\nu(z):=u^\nu(\delta^\nu z).
\]
由步骤 1，只要 $\rho<r/\delta^\nu$，就有
\begin{equation}\label{eq:4-7-5}
  |dv^\nu(0)|=\sup_{B_\rho}|dv^\nu|.
\end{equation}
此外，对每个固定的 $\rho$ 和充分大的 $\nu$，$v^\nu$ 在 $B_\rho$ 中的能量由
$u^\nu$ 在 $B_r$ 中的能量控制。因此，由定理~\ref{thm:4-6-1}，$v^\nu$ 有一个子列，
对某个有限集
\[
  Z=\{z_1,\ldots,z_\ell\}\subset\mathbb C
\]
满足条件 (ii)。式~\eqref{eq:4-7-5} 保证：若 $Z\ne\varnothing$，则 $0\in Z$。
步骤 2 得证。

下面证明，对如此选取的序列 $\psi^\nu,v^\nu$，条件 (iii)、(iv) 成立。特别地，
不必再取子列。下一步包含一个关键恒等式，用来说明极限中没有能量损失。

\medskip
\noindent\textbf{步骤 3.}\quad
\emph{设 $\delta,\delta^\nu$ 如步骤 2，则}
\begin{equation}\label{eq:4-7-6}
  \lim_{R\to\infty}\lim_{\nu\to\infty}
  E(u^\nu;B_{R\delta^\nu})=m_0.
\end{equation}

由假设 (b) 和定理~\ref{thm:4-6-1}，
\begin{equation}\label{eq:4-7-7}
  \lim_{\varepsilon\to0}\lim_{\nu\to\infty}
  E(u^\nu;B_\varepsilon)
  =m_0\ge\hbar>\delta>0.
\end{equation}
接下来构造一个单调不增的实数列 $\varepsilon^\nu\in(0,r]$，使得
\begin{equation}\label{eq:4-7-8}
  \lim_{\nu\to\infty}E(u^\nu;B_{\varepsilon^\nu})=m_0,
  \qquad
  \lim_{\nu\to\infty}\varepsilon^\nu=0.
\end{equation}
细节如下。由式~\eqref{eq:4-7-7}，对每个 $k\in\mathbb N$，都存在
$\varepsilon_k\in(0,r]$，使得
\begin{equation}\label{eq:4-7-9}
  \lim_{\nu\to\infty}
  \bigl|E(u^\nu;B_{\varepsilon_k}(z_0))-m_0\bigr|
  =\left|\lim_{\nu\to\infty}
  E(u^\nu;B_{\varepsilon_k}(z_0))-m_0\right|
  <\frac1k.
\end{equation}
必要时归纳地缩小 $\varepsilon_k$，还可以假设对每个 $k$ 都有
$\varepsilon_{k+1}<\varepsilon_k<1/k$。由式~\eqref{eq:4-7-9}，存在整数列
$\nu_k\in\mathbb N$，使得
\[
  \sup_{\nu\ge\nu_k}
  \bigl|E(u^\nu;B_{\varepsilon_k})-m_0\bigr|<\frac1k,
  \qquad
  \nu_{k+1}>\nu_k.
\]
对 $\nu<\nu_1$ 定义 $\varepsilon^\nu:=r$，并定义
\[
  \varepsilon^\nu:=\varepsilon_k
  \qquad\text{当 }\nu_k\le\nu<\nu_{k+1}.
\]
这个序列满足式~\eqref{eq:4-7-8}。（注意，我们未必能选择 $\varepsilon^\nu$，使
$E(u^\nu;B_{\varepsilon^\nu})$ 等于 $m_0$，因为映射列
$u^\nu:B_r\to M$ 可能由能量小于 $m_0$ 的映射组成，而它们在
$B_r\setminus\{0\}$ 上收敛到常值映射。）我们还有
\begin{equation}\label{eq:4-7-10}
  \lim_{\nu\to\infty}E(u^\nu;B_{R\varepsilon^\nu})=m_0
  \qquad\text{对每个 }R\ge1.
\end{equation}
下极限不小于 $m_0$，由式~\eqref{eq:4-7-8} 和 $R\ge1$ 得出；上极限不大于
$m_0$，则由式~\eqref{eq:4-7-7} 和 $R\varepsilon^\nu\to0$ 得出。

现在反设式~\eqref{eq:4-7-6} 不成立。对 $R>0$ 定义
\[
  f(R):=\limsup_{\nu\to\infty}E(u^\nu;B_{R\delta^\nu}).
\]
这个函数单调不减；而由式~\eqref{eq:4-7-7}，对所有 $R$ 都有 $f(R)\le m_0$。
此外，由步骤 2，函数列
$z\mapsto u^\nu(\delta^\nu z)$ 在 $B_R$ 上\BookRecordIndex{T0247}模冒泡收敛；特别地，当 $R$ 充分大时，
序列 $E(u^\nu;B_{R\delta^\nu})$ 收敛到 $f(R)$。由于式~\eqref{eq:4-7-6}
不成立，便有 $\lim_{R\to\infty}f(R)<m_0$。所以存在常数 $\rho>0$，使得
\begin{equation}\label{eq:4-7-11}
  \limsup_{\nu\to\infty}E(u^\nu;B_{R\delta^\nu})
  \le m_0-\rho
  \qquad\text{对每个 }R>0.
\end{equation}
给定 $R>0$，由式~\eqref{eq:4-7-8} 和式~\eqref{eq:4-7-11}，当 $\nu$ 充分大时，
\[
  E(u^\nu;B_{R\delta^\nu})<E(u^\nu;B_{\varepsilon^\nu}),
\]
因而 $R\delta^\nu<\varepsilon^\nu$。所以 $\delta^\nu/\varepsilon^\nu$ 收敛到零。
此外，由式~\eqref{eq:4-7-4} 和式~\eqref{eq:4-7-10}，
\[
  \lim_{\nu\to\infty}
  E\bigl(u^\nu;A(\delta^\nu,R\varepsilon^\nu)\bigr)
  =\frac\delta2
  \qquad\text{对每个 }R\ge1.
\]
因此，由步骤 2 对 $\delta$ 的选取，当 $\nu$ 充分大时，$u^\nu$ 在\BookRecordIndex{T0170}环域
$A(\delta^\nu,R\varepsilon^\nu)$ 上的限制满足引理~\ref{lem:4-7-3} 的假设，
其中 $\mu$ 是步骤 2 中固定的常数。取 $T:=\log R$，便得到
\[
  \begin{aligned}
  \limsup_{\nu\to\infty}
  E\bigl(u^\nu;A(R\delta^\nu,\varepsilon^\nu)\bigr)
  &\le\frac{4^\mu}{R^{2\mu}}\lim_{\nu\to\infty}
  E\bigl(u^\nu;A(\delta^\nu,R\varepsilon^\nu)\bigr)\\
  &=\frac{2^{2\mu-1}\delta}{R^{2\mu}}
  \qquad\text{对每个 }R\ge2.
  \end{aligned}
\]
于是，由式~\eqref{eq:4-7-11}，
\[
  \lim_{\nu\to\infty}E(u^\nu;B_{\varepsilon^\nu})
  \le m_0-\rho+\frac{2^{2\mu-1}\delta}{R^{2\mu}}
  \qquad\text{对每个 }R\ge2.
\]
取充分大的 $R$，使最后一项小于 $\rho$，这与式~\eqref{eq:4-7-8} 矛盾。
步骤 3 得证。

\medskip
\noindent\textbf{步骤 4.}\quad
\emph{完成证明。}

先证明 (iii)。由式~\eqref{eq:4-7-6} 以及步骤 2 中 $\delta^\nu,v^\nu$ 的定义，
\[
  \lim_{R\to\infty}\lim_{\nu\to\infty}E(v^\nu;B_R)
  =\lim_{R\to\infty}\lim_{\nu\to\infty}
  E(u^\nu;B_{R\delta^\nu})
  =m_0;
\]
并且对每个 $\nu$，
\[
  E(v^\nu;B_1)=E(u^\nu;B_{\delta^\nu})
  =m_0-\frac\delta2\ge m_0-\frac\hbar2.
\]
所以，映射列 $v^\nu$ 的所有\BookRecordIndex{T0275}气泡点 $z_1,\ldots,z_\ell$ 都位于闭单位圆盘
$B_1$ 中。现固定 $s>1$。再次使用式~\eqref{eq:4-7-6}，得到
\[
  \begin{aligned}
  m_0
  &=\lim_{R\to\infty}\lim_{\nu\to\infty}E(v^\nu;B_R)\\
  &=\lim_{R\to\infty}\lim_{\nu\to\infty}
    E(v^\nu;B_R\setminus B_s)
    +\lim_{\nu\to\infty}E(v^\nu;B_s)\\
  &=\lim_{R\to\infty}E(v;B_R\setminus B_s)
    +\lim_{\nu\to\infty}E(v^\nu;B_s)\\
  &=E(v;\mathbb C\setminus B_s)
    +\lim_{\varepsilon\to0}\lim_{\nu\to\infty}
    E\left(v^\nu;
      B_s\setminus\bigcup_{j=1}^{\ell}B_\varepsilon(z_j)\right)
    +\sum_{j=1}^{\ell}m_j\\
  &=E(v)+\sum_{j=1}^{\ell}m_j.
  \end{aligned}
\]
这证明了 (iii)。

再证明 (iv)。我们已经看到，$v^\nu$ 的所有\BookRecordIndex{T0275}气泡点都位于闭单位圆盘内。因此，若 $v$
是常值映射且 $1<R_1<R_2$，则
\[
  \lim_{\nu\to\infty}E\bigl(v^\nu;A(R_1,R_2)\bigr)
  =E\bigl(v;A(R_1,R_2)\bigr)=0.
\]
这说明，对 $R>1$，$E(v^\nu;B_R)=E(u^\nu;B_{R\delta^\nu})$ 在
$\nu\to\infty$ 时的极限与 $R$ 无关。因此，由式~\eqref{eq:4-7-6}，
\[
  \lim_{\nu\to\infty}E(v^\nu;B_R)=m_0
  \qquad\text{对每个 }R>1.
\]
由于
\[
  E(v^\nu;B_1)=m_0-\delta/2,
\]
这只有在单位圆盘的边界上发生\BookRecordIndex{T0243}冒泡时才可能。也就是说，$v^\nu$ 的奇异集 $Z$ 包含某个
满足 $|z_j|=1$ 的点 $z_j$。但步骤 2 已说明：若 $Z\ne\varnothing$，则 $0\in Z$。
所以如所要求的，$\ell=\#Z\ge2$。这证明了命题~\ref{prop:4-7-1}。
\end{proof}
为了完成本节的证明，还须证明引理~\ref{lem:4-7-3}。

\begin{proof}[引理~\ref{lem:4-7-3} 的证明]
选取正数 $\delta,C$，使引理~\ref{lem:4-3-1}(i) 的结论在 $c$ 换成 $2C$ 时成立。
于是，对每个 $r>0$ 和每条 \BookRecordIndex{T0003}$J$-全纯曲线 $v:B_r\to M$，都有
\begin{equation}\label{eq:4-7-12}
  E(v;B_r)<\delta
  \quad\Longrightarrow\quad
  \frac12|dv(0)|^2\le\frac C{r^2}E(v;B_r).
\end{equation}
必要时缩小 $\delta$，并选取适当的常数 $0<\mu<1$，还可以假设
定理~\ref{thm:4-4-1}(iii) 的结论在 $\delta$ 换成
$3\pi\sqrt{C\delta}$、$c$ 换成 $1/(4\pi\mu)$ 时成立。也就是说，对每条光滑闭路
$\gamma:S^1\to M$，有
\begin{equation}\label{eq:4-7-13}
  \ell(\gamma)<3\pi\sqrt{C\delta}
  \quad\Longrightarrow\quad
  |a(\gamma)|\le\frac{\ell(\gamma)^2}{4\pi\mu}.
\end{equation}
定义
\begin{equation}\label{eq:4-7-14}
  c:=16\sqrt C\left(2\pi+\frac1\mu\right).
\end{equation}
我们将证明，引理~\ref{lem:4-7-3} 对这些常数 $\delta,c$ 成立。

首先，不失一般性，可以假设 $r=1/R$。事实上，定义
\[
  \lambda:=\sqrt{rR},
\]
并把 $u,r,R$ 分别替换为
\[
  \widetilde u(z):=u(\lambda z),
  \qquad
  \widetilde r:=r/\lambda,
  \qquad
  \widetilde R:=R/\lambda.
\]
于是 $\widetilde r\widetilde R=1$；若引理已对三元组
$(\widetilde u,\widetilde r,\widetilde R)$ 成立，它也对 $(u,r,R)$ 成立。

令 $R>0$，并设 $u:A(1/R,R)\to M$ 是能量满足
\[
  E(u)<\delta
\]
的 \BookRecordIndex{T0003}$J$-全纯曲线。定义 \BookRecordIndex{T0003}$J$-全纯曲线
\[
  v:[-\log R,\log R]\times\mathbb R/2\pi\mathbb Z\to M
\]
为
\[
  v(s,t):=u(e^{s+it}),
  \qquad
  -\log R<s<\log R,
  \qquad
  t\in\mathbb R/2\pi\mathbb Z.
\]
对 $s\in[-\log R,\log R]$，定义
$\gamma_s:\mathbb R/2\pi\mathbb Z\to M$ 为
\[
  \gamma_s(t):=u(e^{s+it})=v(s,t).
\]
由于 $E(v)=E(u)<\delta$，先验估计~\eqref{eq:4-7-12} 在 $r=\log2$ 时给出
\begin{equation}\label{eq:4-7-15}
  |\dot\gamma_s(t)|^2=|\partial_sv(s,t)|^2
  \le\frac C{(\log2)^2}E(u),
  \qquad |s|\le\log(R/2).
\end{equation}
由此
\[
  \begin{aligned}
  \ell(\gamma_s)^2
  &=\left(\int_0^{2\pi}|\dot\gamma_s(t)|\,dt\right)^2\\
  &\le2\pi\int_0^{2\pi}|\dot\gamma_s(t)|^2\,dt\\
  &\le\frac{4\pi^2C}{(\log2)^2}E(u)\\
  &\le9\pi^2CE(u).
  \end{aligned}
\]
这里使用了不等式 $4/(\log2)^2\le9$。

先证明式~\eqref{eq:4-7-1}。回忆注~\ref{rem:4-4-3}：充分短的闭路
$\gamma:S^1\to M$ 的\BookRecordIndex{T0202}局部辛作用量定义为
\[
  a(\gamma):=-\int_{B_1}u_\gamma^*\omega,
\]
其中 $u_\gamma:B_1\to M$ 是 $\gamma$ 的一个局部延拓。由于当
$|s|\le\log(R/2)$ 时，
\[
  \ell(\gamma_s)\le3\pi\sqrt{CE(u)}
  \le3\pi\sqrt{C\delta},
\]
由式~\eqref{eq:4-7-13} 可得
\begin{equation}\label{eq:4-7-16}
  |a(\gamma_s)|\le\frac{\ell(\gamma_s)^2}{4\pi\mu},
  \qquad |s|\le\log(R/2).
\end{equation}
此外，当
$-\log(R/2)\le s_0\le s_1\le\log(R/2)$ 时，
\begin{equation}\label{eq:4-7-17}
  E\bigl(v;[s_0,s_1]\times\mathbb R/2\pi\mathbb Z\bigr)
  +a(\gamma_{s_1})-a(\gamma_{s_0})=0.
\end{equation}
要看出这一点，用小圆盘封住圆柱
$v:[s_0,s_1]\times\mathbb R/2\pi\mathbb Z\to M$ 的两端，得到一个可缩球面，
再在这个球面上积分\BookRecordIndex{T0363}辛形式。

现在考虑函数
\[
  \varepsilon(T):=E\bigl(u;A(e^T/R,e^{-T}R)\bigr),
  \qquad \log2\le T\le\log R.
\]
简记 $\rho:=\log R$，则
\[
  \begin{aligned}
  \varepsilon(T)
  &=E\bigl(v;[-\rho+T,\rho-T]\times\mathbb R/2\pi\mathbb Z\bigr)\\
  &=\int_{-\rho+T}^{\rho-T}\int_0^{2\pi}
    |\partial_sv(s,t)|^2\,dt\,ds.
  \end{aligned}
\]
因此，由式~\eqref{eq:4-7-16} 和式~\eqref{eq:4-7-17}，
\[
  \begin{aligned}
  \varepsilon(T)
  &=a(\gamma_{-\rho+T})-a(\gamma_{\rho-T})\\
  &\le\frac1{4\pi\mu}
    \left(\ell(\gamma_{-\rho+T})^2
    +\ell(\gamma_{\rho-T})^2\right)\\
  &\le\frac1{2\mu}\int_0^{2\pi}
    |\partial_sv(-\rho+T,t)|^2\,dt
    +\frac1{2\mu}\int_0^{2\pi}
    |\partial_sv(\rho-T,t)|^2\,dt\\
  &=-\frac1{2\mu}\dot\varepsilon(T).
  \end{aligned}
\]
这蕴含
\[
  \dot\varepsilon(T)\le-2\mu\varepsilon(T)<0,
\]
从而
\begin{equation}\label{eq:4-7-18}
  \varepsilon(T)
  \le e^{-2\mu(T-\log2)}\varepsilon(\log2)
  \le4^\mu e^{-2\mu T}E(u),
  \qquad \log2\le T\le\log R.
\end{equation}
式~\eqref{eq:4-7-1} 得证。

再证明式~\eqref{eq:4-7-2}。第一步是建立估计
\begin{equation}\label{eq:4-7-19}
  |\partial_sv(s,t)|^2
  \le\frac{C4^\mu e^{2\mu}}{(\log2)^2}
  e^{-2\mu(\rho-|s|)}E(v),
  \qquad |s|\le\rho-\log2.
\end{equation}
为此，先假设
\[
  \rho-\log2-1\le|s|\le\rho-\log2.
\]
由式~\eqref{eq:4-7-15}，
\[
  |\partial_sv(s,t)|^2
  \le\frac C{(\log2)^2}E(v)
  \le\frac{C4^\mu e^{2\mu}}{(\log2)^2}
  e^{-2\mu(\rho-|s|)}E(v).
\]
最后一个不等式使用了
$\rho-|s|\le\log2+1$，因而
$e^{2\mu(\rho-|s|)}\le4^\mu e^{2\mu}$。现在假设
\[
  |s|\le\rho-\log2-1,
\]
并记 $T:=\rho-|s|$，则 $T-1\ge\log2$。于是
\[
  \begin{aligned}
  |\partial_sv(s,t)|^2
  &\le C E\bigl(v;
    [-\rho+T-1,\rho-T+1]\times\mathbb R/2\pi\mathbb Z\bigr)\\
  &=C\varepsilon(T-1)\\
  &\le C4^\mu e^{2\mu}e^{-2\mu T}E(v).
  \end{aligned}
\]
这里第一个不等式由式~\eqref{eq:4-7-12} 在 $r=1$ 时给出，最后一个不等式由
式~\eqref{eq:4-7-18} 中把 $T$ 换成 $T-1$ 得出。由于
$T=\rho-|s|$ 且 $(\log2)^2<1$，这就蕴含式~\eqref{eq:4-7-19}。

当 $s\ge0$ 时，估计~\eqref{eq:4-7-19} 蕴含
\[
  0\le s\le\rho-\log2
  \quad\Longrightarrow\quad
  |\partial_sv(s,t)|
  \le8\sqrt C e^{-\mu(\rho-s)}\sqrt{E(v)}.
\]
这里使用了不等式
$2^\mu e^\mu/\log2\le2e/\log2<8$。现固定一点
$(s_0,t_0)\in\mathbb R^2$，使得
$0\le s_0\le\rho-T$ 且 $0\le t_0\le2\pi$。对最后一个不等式积分，得到
\[
  \begin{aligned}
  d\bigl(v(0,0),v(s_0,t_0)\bigr)
  &\le\int_0^{s_0}|\partial_sv(s,0)|\,ds
    +\int_0^{t_0}|\partial_tv(s_0,t)|\,dt\\
  &\le8\sqrt C\left(
    \int_0^{s_0}e^{-\mu(\rho-s)}\,ds
    +\int_0^{t_0}e^{-\mu(\rho-s_0)}\,dt\right)\sqrt{E(v)}\\
  &\le8\sqrt C\left(
    \int_{-\infty}^{-T}e^{\mu s}\,ds
    +2\pi e^{-\mu T}\right)\sqrt{E(v)}\\
  &=8\sqrt C\left(\frac1\mu+2\pi\right)
    e^{-\mu T}\sqrt{E(v)}.
  \end{aligned}
\]
利用关于 $s=0$ 的对称性，可得
\[
  d\bigl(v(s_0,t_0),v(s_1,t_1)\bigr)
  \le16\sqrt C\left(\frac1\mu+2\pi\right)
  e^{-\mu T}\sqrt{E(v)}
\]
对 $\log2\le T\le\rho$、
$s_0,s_1\in[-\rho+T,\rho-T]$ 以及 $t_0,t_1\in[0,2\pi]$ 成立。
因此，式~\eqref{eq:4-7-2} 对式~\eqref{eq:4-7-14} 给出的 $c$ 成立。
引理~\ref{lem:4-7-3} 得证。
\end{proof}

\begin{exercise}\label{ex:4-7-6}
把本节的结果推广到具有拉格朗日边界条件的 \BookRecordIndex{T0003}$J$-全纯曲线（参见 \srcite{120}）。
\end{exercise}
